\documentclass[10pt,a4paper]{amsart}
\usepackage{amsmath,amssymb,amsthm,graphicx}
% Retained minimal preamble for the Stacks Project chapter sites-cohomology.tex (Cohomology on Sites).
% Verbatim excerpt of the chapter preamble: Xy-pic setup (native preamble lines 14-18) and the native
% theorem declarations with the mathematical macros (native preamble lines 154-217). Copied mechanically;
% nothing added, removed or redefined. The wrapper class and the amsmath/amssymb/amsthm/graphicx base
% packages are supplied by the assembler, so the class/externaldocument part of the native preamble is omitted.
% For commutative diagrams we use Xy-pic
\usepackage[all]{xy}

% We use 2cell for 2-commutative diagrams.
\xyoption{2cell}
% Theorem environments.
%
\theoremstyle{plain}
\newtheorem{theorem}[subsection]{Theorem}
\newtheorem{proposition}[subsection]{Proposition}
\newtheorem{lemma}[subsection]{Lemma}

\theoremstyle{definition}
\newtheorem{definition}[subsection]{Definition}
\newtheorem{example}[subsection]{Example}
\newtheorem{exercise}[subsection]{Exercise}
\newtheorem{situation}[subsection]{Situation}

\theoremstyle{remark}
\newtheorem{remark}[subsection]{Remark}
\newtheorem{remarks}[subsection]{Remarks}

\numberwithin{equation}{subsection}

% Macros
%
\def\lim{\mathop{\mathrm{lim}}\nolimits}
\def\colim{\mathop{\mathrm{colim}}\nolimits}
\def\Spec{\mathop{\mathrm{Spec}}}
\def\Hom{\mathop{\mathrm{Hom}}\nolimits}
\def\Ext{\mathop{\mathrm{Ext}}\nolimits}
\def\SheafHom{\mathop{\mathcal{H}\!\mathit{om}}\nolimits}
\def\SheafExt{\mathop{\mathcal{E}\!\mathit{xt}}\nolimits}
\def\Sch{\mathit{Sch}}
\def\Mor{\mathop{\mathrm{Mor}}\nolimits}
\def\Ob{\mathop{\mathrm{Ob}}\nolimits}
\def\Sh{\mathop{\mathit{Sh}}\nolimits}
\def\NL{\mathop{N\!L}\nolimits}
\def\CH{\mathop{\mathrm{CH}}\nolimits}
\def\proetale{{pro\text{-}\acute{e}tale}}
\def\etale{{\acute{e}tale}}
\def\QCoh{\mathit{QCoh}}
\def\Ker{\mathop{\mathrm{Ker}}}
\def\Im{\mathop{\mathrm{Im}}}
\def\Coker{\mathop{\mathrm{Coker}}}
\def\Coim{\mathop{\mathrm{Coim}}}

% Boxtimes
%
\DeclareMathSymbol{\boxtimes}{\mathbin}{AMSa}{"02}

%
% Macros for moduli stacks/spaces
%
\def\QCohstack{\mathcal{QC}\!\mathit{oh}}
\def\Cohstack{\mathcal{C}\!\mathit{oh}}
\def\Spacesstack{\mathcal{S}\!\mathit{paces}}
\def\Quotfunctor{\mathrm{Quot}}
\def\Hilbfunctor{\mathrm{Hilb}}
\def\Curvesstack{\mathcal{C}\!\mathit{urves}}
\def\Polarizedstack{\mathcal{P}\!\mathit{olarized}}
\def\Complexesstack{\mathcal{C}\!\mathit{omplexes}}
% \Pic is the operator that assigns to X its picard group, usage \Pic(X)
% \Picardstack_{X/B} denotes the Picard stack of X over B
% \Picardfunctor_{X/B} denotes the Picard functor of X over B
\def\Pic{\mathop{\mathrm{Pic}}\nolimits}
\def\Picardstack{\mathcal{P}\!\mathit{ic}}
\def\Picardfunctor{\mathrm{Pic}}
\def\Deformationcategory{\mathcal{D}\!\mathit{ef}}
\usepackage{cleveref}
\crefname{originalreference}{original reference}{original references}
\Crefname{originalreference}{Original reference}{Original references}
\makeatletter
\newlabel{lemma-collect-true-things-Zar}{{0DCU}{source}{}{sourcepointer-1}{}}
\newlabel{lemma-collect-true-things-Zar@cref}{{[originalreference][0][]0DCU}{[1][0][]source}{}{}{}}
\makeatother
\crefname{originalreference}{original reference}{original references}
\Crefname{originalreference}{Original reference}{Original references}
\makeatletter
\newlabel{lemma-describe-pullback-pi}{{09X3}{source}{}{sourcepointer-2}{}}
\newlabel{lemma-describe-pullback-pi@cref}{{[originalreference][0][]09X3}{[1][0][]source}{}{}{}}
\makeatother
\crefname{originallemma}{original lemma}{original lemmas}
\Crefname{originallemma}{Original lemma}{Original lemmas}
\makeatletter
\newlabel{sites-lemma-composition-morphisms-sites}{{03CB}{source}{}{sourcepointer-3}{}}
\newlabel{sites-lemma-composition-morphisms-sites@cref}{{[originallemma][0][]03CB}{[1][0][]source}{}{}{}}
\makeatother
\crefname{originallemma}{original lemma}{original lemmas}
\Crefname{originallemma}{Original lemma}{Original lemmas}
\makeatletter
\newlabel{sites-lemma-relocalize}{{03EH}{source}{}{sourcepointer-4}{}}
\newlabel{sites-lemma-relocalize@cref}{{[originallemma][0][]03EH}{[1][0][]source}{}{}{}}
\makeatother
\crefname{originallemma}{original lemma}{original lemmas}
\Crefname{originallemma}{Original lemma}{Original lemmas}
\makeatletter
\newlabel{sites-lemma-relocalize-given-fibre-products}{{09W9}{source}{}{sourcepointer-5}{}}
\newlabel{sites-lemma-relocalize-given-fibre-products@cref}{{[originallemma][0][]09W9}{[1][0][]source}{}{}{}}
\makeatother
\title{The comparison diagrams of the qc and Zariski topologies on locally quasi-compact spaces commute with a proper morphism}
\author{The Stacks Project Authors}
\date{stacks-project:sites-cohomology.tex:local-snapshot-20261011; GNU FDL 1.2 or later}
\begin{document}
\maketitle
\noindent\textit{Editorial note.} All delivered mathematical text is the author's own, copied from the pinned Stacks Project chapter file sites-cohomology.tex; no statement or proof step was rewritten, reordered or removed.\par\medskip
\noindent\textit{Editorial note.} The counted proof is the author's complete argument: the proof is written out in full: the identification of f\_qc and the commutativity of the first diagram, the derivation of the second from the collected Zariski properties, and the reduction of the push-pull identity to part (7a) of those properties.\par\medskip
\noindent\textit{Editorial note.} The label lemma-push-pull-LC names the counted statement, and the proof environment immediately following it in the source is delivered verbatim as its proof.\par\medskip
\noindent\textit{Editorial note.} The counted statement is carried by the author's own lemma environment in the chapter source.\par\medskip
\noindent\textit{Editorial note.} The statement and proof are copied from the pinned chapter file; the five cited inputs are kept as external pointers with their canonical Stacks tags.\par\medskip
\noindent\textit{Editorial note.} The preamble is the chapter's own retained minimal preamble (sources/stacks-sites-cohomology/retained-preamble.tex), a verbatim excerpt of the native preamble: the Xy-pic setup and the native theorem declarations with the mathematical macros.\par\medskip
\noindent\textit{Allowed standard prerequisites (editorial).}\begin{itemize}
\item The morphisms of topoi f\_qc, epsilon, pi and the comparison a = pi o epsilon for LC\_qc. Same chapter, the construction of the morphisms of topoi between the qc, Zariski and topological realisations of a morphism of locally quasi-compact spaces. Supplies every arrow of the two displayed diagrams.
\item Relocalisation and the fibre-product description of the induced morphism. Sites, the lemmas describing the morphism of topoi induced by relocalisation and the continuous functor Z/Y |-> Z x\_Y X/X. Supplies the identification of f\_qc used in the proof.
\item Collected properties of the Zariski comparison. Same chapter, the lemma collecting the compatibility statements for the Zariski topology, in particular parts (5) and (7a). Supplies the commutativity of the second diagram and the final identification pi\_Y\textasciicircum{}\{-1\}f\_*F = f\_\{Zar,*\}pi\_X\textasciicircum{}\{-1\}F.
\end{itemize}\medskip
\ifdefined\hypertarget\hypertarget{sourcepointer-1}{}\fi
\noindent\textit{Source reference.} Original reference 0DCU: sources/stacks-sites-cohomology/pinned-original-source.tex, source lines 7770--7770. Prerequisite pointer invoked at proof line(s) 7962; it names a standard fact of the same chapter that the counted argument uses, and every step of the counted proof built on it is retained verbatim in the delivered body.\par\medskip
\ifdefined\hypertarget\hypertarget{sourcepointer-2}{}\fi
\noindent\textit{Source reference.} Original reference 09X3: sources/stacks-sites-cohomology/pinned-original-source.tex, source lines 7675--7675. Prerequisite pointer invoked at proof line(s) 7969; it names a standard fact of the same chapter that the counted argument uses, and every step of the counted proof built on it is retained verbatim in the delivered body.\par\medskip
\ifdefined\hypertarget\hypertarget{sourcepointer-3}{}\fi
\noindent\textit{Source reference.} Original lemma 03CB: sources/stacks-chow/tags.txt, source lines 4378--4378. Prerequisite pointer invoked at proof line(s) 7960; it names a standard fact of a sibling chapter that the counted argument uses, and every step of the counted proof built on it is retained verbatim in the delivered body.\par\medskip
\ifdefined\hypertarget\hypertarget{sourcepointer-4}{}\fi
\noindent\textit{Source reference.} Original lemma 03EH: sources/stacks-chow/tags.txt, source lines 4456--4456. Prerequisite pointer invoked at proof line(s) 7954; it names a standard fact of a sibling chapter that the counted argument uses, and every step of the counted proof built on it is retained verbatim in the delivered body.\par\medskip
\ifdefined\hypertarget\hypertarget{sourcepointer-5}{}\fi
\noindent\textit{Source reference.} Original lemma 09W9: sources/stacks-chow/tags.txt, source lines 12487--12487. Prerequisite pointer invoked at proof line(s) 7957; it names a standard fact of a sibling chapter that the counted argument uses, and every step of the counted proof built on it is retained verbatim in the delivered body.\par\medskip
\begin{lemma}
\label{lemma-push-pull-LC}
Let $f : X \to Y$ be a morphism of $\textit{LC}_{qc}$.
Then there are commutative diagrams of topoi
$$
\vcenter{
\xymatrix{
\Sh(\textit{LC}_{qc}/X) \ar[r]_{f_{qc}} \ar[d]_{\epsilon_X} &
\Sh(\textit{LC}_{qc}/Y) \ar[d]^{\epsilon_Y} \\
\Sh(\textit{LC}_{Zar}/X) \ar[r]^{f_{Zar}} &
\Sh(\textit{LC}_{Zar}/Y)
}
}
\quad\text{and}\quad
\vcenter{
\xymatrix{
\Sh(\textit{LC}_{qc}/X) \ar[r]_{f_{qc}} \ar[d]_{a_X} &
\Sh(\textit{LC}_{qc}/Y) \ar[d]^{a_Y} \\
\Sh(X) \ar[r]^f &
\Sh(Y)
}
}
$$
with $a_X = \pi_X \circ \epsilon_X$, $a_Y = \pi_X \circ \epsilon_X$.
If $f$ is proper, then $a_Y^{-1} \circ f_* = f_{qc, *} \circ a_X^{-1}$.
\end{lemma}

\begin{proof}
The morphism of topoi $f_{qc}$ is the one from
Sites, Lemma \ref{sites-lemma-relocalize}
which in our case comes from the continuous functor
$Z/Y \mapsto Z \times_Y X/X$, see
Sites, Lemma \ref{sites-lemma-relocalize-given-fibre-products}.
The diagram on the left commutes because the corresponding
continuous functors compose correctly
(see Sites, Lemma \ref{sites-lemma-composition-morphisms-sites}).
The diagram on the right commutes because the one on the left does
and because of part (5) of Lemma \ref{lemma-collect-true-things-Zar}.

\medskip\noindent
Proof of the final assertion. The reader may repeat the proof of part (7a) of
Lemma \ref{lemma-collect-true-things-Zar}; we will instead deduce this from it.
As $\epsilon_{Y, *}$ is the identity functor on underlying presheaves,
it reflects isomorphisms. The description
in Lemma \ref{lemma-describe-pullback-pi}
shows that $\epsilon_{Y, *} \circ a_Y^{-1} = \pi_Y^{-1}$
and similarly for $X$. To show that the canonical map
$a_Y^{-1}f_*\mathcal{F} \to f_{qc, *}a_X^{-1}\mathcal{F}$
is an isomorphism, it suffices to show that
$$
\pi_Y^{-1}f_*\mathcal{F} =
\epsilon_{Y, *}a_Y^{-1}f_*\mathcal{F} \to
\epsilon_{Y, *}f_{qc, *}a_X^{-1}\mathcal{F} =
f_{Zar, *}\epsilon_{X, *}a_X^{-1}\mathcal{F} =
f_{Zar, *}\pi_X^{-1}\mathcal{F}
$$
is an isomorphism. This is part
(7a) of Lemma \ref{lemma-collect-true-things-Zar}.
\end{proof}

\end{document}
